% Created 2026-08-16 Sun 19:38
% Intended LaTeX compiler: pdflatex
\documentclass[babel]{beamer}

\usepackage[hyperref]{biblatex}
\addbibresource{~/Documents/GLOBAL_ORG/org/roam/bib/master.bib}
\institute[]{Università degli Studi di Torino}
\newcommand{\use}[2][]{\usepackage[#1]{#2}}
% PACCHETTI FONDAMENTLAI
\use[utf8]{inputenc}
\use[T1]{fontenc}
\use{graphicx}
\use{longtable}
\use{wrapfig}
\use{rotating}
\use[normalem]{ulem}
\use{amsmath}
\use{amsthm}
\use{amssymb}
\use{capt-of}
\use[italian]{babel}
\use[babel]{csquotes}
\use{microtype}
\use{lmodern}
\use{subfig} % sottofigure
\use{multicol} % due colonne
\use{lipsum} % lorem ipsum
\use{color} % colori in latex
\use{parskip} % rimuove l'indentazione dei nuovi paragrafi %% Add parbox=false to all new tcolorbox
\use{centernot}
\use[outline]{contour}\contourlength{3pt}
\use{fancyhdr}
\use{layout}
\use[most]{tcolorbox} % Riquadri colorati
\use{ifthen} % IFTHEN
\use{geometry}

% pacchetti matematica
\use{yhmath}
\use{dsfont}
\use{mathrsfs}
\use{cancel} % semplificare
\use{polynom} %divisione tra polinomi
\use{forest} % grafi ad albero
\use{booktabs} % tabelle
\use{commath} %simboli e differenziali
\use{bm} %bold
\use[fulladjust]{marginnote} %to use marginnote for date notes
\use{arrayjobx}%array
\use[intlimits]{empheq} % Riquadri colorati attorno alle equazioni
\use{mathtools}
\use{circuitikz} % Disegnare i circuiti
%%%%%%%%%%%%%


%%%% QUIVER
\newcommand{\duepunti}{\,\mathchar\numexpr"6000+`:\relax\,}
% A TikZ style for curved arrows of a fixed height, due to AndréC.
\tikzset{curve/.style={settings={#1},to path={(\tikztostart)
    .. controls ($(\tikztostart)!\pv{pos}!(\tikztotarget)!\pv{height}!270:(\tikztotarget)$)
    and ($(\tikztostart)!1-\pv{pos}!(\tikztotarget)!\pv{height}!270:(\tikztotarget)$)
    .. (\tikztotarget)\tikztonodes}},
    settings/.code={\tikzset{quiver/.cd,#1}
        \def\pv##1{\pgfkeysvalueof{/tikz/quiver/##1}}},
    quiver/.cd,pos/.initial=0.35,height/.initial=0}

% TikZ arrowhead/tail styles.
\tikzset{tail reversed/.code={\pgfsetarrowsstart{tikzcd to}}}
\tikzset{2tail/.code={\pgfsetarrowsstart{Implies[reversed]}}}
\tikzset{2tail reversed/.code={\pgfsetarrowsstart{Implies}}}
% TikZ arrow styles.
\tikzset{no body/.style={/tikz/dash pattern=on 0 off 1mm}}
%%%%%%%%%%


%% DEFINIZIONI COMANDI MATEMATICI
\let\sin\relax %TOGLIE LA DEFINIZIONE SU "\sin"

% cambia la definizione di empty set
% ---
\let\oldemptyset\emptyset
% ---
% \let\emptyset\varnothing
% ---
% \let\emptyset\relax
% \newcommand{\emptyset}{\text{\textnormal{\O}}}
% ---

\DeclareMathOperator{\bounded}{bd}
\DeclareMathOperator{\sin}{sen}
\DeclareMathOperator{\epi}{Epi}
\DeclareMathOperator{\cl}{cl}
\DeclareMathOperator{\graph}{graph}
\DeclareMathOperator{\arcsec}{arcsec}
\DeclareMathOperator{\arccot}{arccot}
\DeclareMathOperator{\arccsc}{arccsc}
\DeclareMathOperator{\spettro}{Spettro}
\DeclareMathOperator{\nulls}{nullspace}
\DeclareMathOperator{\dom}{dom}
\DeclareMathOperator{\ar}{ar}
\DeclareMathOperator{\const}{Const}
\DeclareMathOperator{\fun}{Fun}
\DeclareMathOperator{\rel}{Rel}
\DeclareMathOperator{\altezza}{ht}
\let\det\relax %TOGLIE LA DEFINIZIONE SU "\det"
\DeclareMathOperator{\det}{det}
\DeclareMathOperator{\End}{End}
\DeclareMathOperator{\gl}{GL}
\DeclareMathOperator{\Id}{Id}
\DeclareMathOperator{\id}{Id}
\DeclareMathOperator{\I}{\mathds{1}}
\DeclareMathOperator{\II}{II}
\DeclareMathOperator{\rank}{rank}
\DeclareMathOperator{\tr}{tr}
\DeclareMathOperator{\tc}{t.c.}
\DeclareMathOperator{\T}{T}
\DeclareMathOperator{\var}{Var}
\DeclareMathOperator{\cov}{Cov}
\DeclareMathOperator{\st}{st}
\DeclareMathOperator{\mon}{Mon}
\newcommand{\card}[1]{\left\vert #1 \right\vert}
\newcommand{\trasposta}[1]{\prescript{\text{T}}{}{#1}}
\newcommand{\1}{\mathds{1}}
\newcommand{\R}{\mathds{R}}
\newcommand{\diesis}{\#}
\newcommand{\bemolle}{\flat}
\newcommand{\starR}{\nonstandard{\R}}
\newcommand{\borel}{\mathscr{B}}
\newcommand{\lebesgue}[1]{\mathscr{L}\left(#1\right)}
\newcommand{\media}{\mathds{E}}
\newcommand{\K}{\mathds{K}}
\newcommand{\A}{\mathds{A}}
\newcommand{\Q}{\mathds{Q}}
\newcommand{\N}{\mathds{N}}
\newcommand{\C}{\mathds{C}}
\newcommand{\Z}{\mathds{Z}}
\newcommand{\qo}{\hspace{1em}\text{q.o.}\,}
\renewcommand{\tilde}[1]{\widetilde{#1}}
\renewcommand{\parallel}{\mathrel{/\mkern-5mu/}}
\newcommand{\parti}[2][]{\wp_{#1}(#2)}
\newcommand{\diff}[1]{\operatorname{d}_{#1}}
\let\oldvec\vec
\renewcommand{\vec}[1]{\overrightarrow{\vphantom{i}#1}}
\newcommand{\floor}[1]{\left\lfloor #1 \right\rfloor}
\newcommand{\cat}[1]{\mathbf{#1}}
\newcommand{\dfreccia}[1]{\xrightarrow{\ #1 \ }}
\newcommand{\sfreccia}[1]{\xleftarrow{\ #1 \ }}
\newcommand{\formalsum}[2]{{\sum_{#1}^{#2}}{\vphantom{\sum}}'}
\newcommand{\minim}[2]{\mu_{#1}\, \left(#2\right)}
\newcommand{\concat}{\null^{\frown}} % concatenazione di stringe

%% Definizione di \dotminus

\makeatletter
\newcommand{\dotminus}{\mathbin{\text{\@dotminus}}}

\newcommand{\@dotminus}{%
  \ooalign{\hidewidth\raise1ex\hbox{.}\hidewidth\cr$\m@th-$\cr}%
}
\makeatother

%tramite i prossimi due comandi posso decidere come scrivere i logaritmi naturali in tutti i documenti: ho infatti eliminato qualsiasi differenza tra "ln" e "log": se si vuole qualcosa di diverso bisogna inserire manualmente il tutto
\let\ln\relax
\DeclareMathOperator{\ln}{ln}
\let\log\relax
\DeclareMathOperator{\log}{log}
%%%%%%

%% NUOVI COMANDI
\newcommand{\straniero}[1]{\textit{#1}} %parole straniere
\newcommand{\titolo}[1]{\textsc{#1}} %titoli
\newcommand{\qedd}{\tag*{$\blacksquare$}} %qed per ambienti matemastici
\renewcommand{\qedsymbol}{$\blacksquare$} %modifica colore qed
\newcommand{\ooverline}[1]{\overline{\overline{#1}}}
\newcommand{\circoletto}[1]{\left(#1\right)^{\text{o}}}
%
\newcommand{\qmatrice}[1]{\begin{pmatrix}
#1_{11} & \cdots & #1_{1n}\\
\vdots & \ddots & \vdots \\
#1_{m1} & \cdots & #1_{mn}
\end{pmatrix}}
%
\newcommand{\parentesi}[2]{%
\underset{#1}{\underbrace{#2}}%
}
%
\newcommand{\norma}[1]{% Norma
\left\lVert#1\right\rVert%
}
\newcommand{\scalare}[2]{% Scalare
\left\langle #1, #2\right\rangle
}
%%%%%

%%%% Change footnote appearance
%%%%

\makeatletter
% ---- marker nel TESTO: (nota 1)
\renewcommand\@makefnmark{%
  \hbox{\normalfont\footnotesize(nota~\@thefnmark)}%
}

% ---- layout e marker diverso in PIÉ DI PAGINA: (1)
\renewcommand\@makefntext[1]{%
  \parindent 1em
  \noindent
  % qui non chiamo \@makefnmark, ma uso direttamente (\@thefnmark)
  \hb@xt@1.8em{\hss\normalfont\footnotesize(\@thefnmark)} %
  #1%
}
\makeatother
%%%%
%%%%

%% RESTRIZIONI
\newcommand{\referenze}[2]{
	\phantomsection{}#2\textsuperscript{\textcolor{blue}{\textbf{#1}}}
}

\let\restriction\relax

\def\restriction#1#2{\mathchoice
              {\setbox1\hbox{${\displaystyle #1}_{\scriptstyle #2}$}
              \restrictionaux{#1}{#2}}
              {\setbox1\hbox{${\textstyle #1}_{\scriptstyle #2}$}
              \restrictionaux{#1}{#2}}
              {\setbox1\hbox{${\scriptstyle #1}_{\scriptscriptstyle #2}$}
              \restrictionaux{#1}{#2}}
              {\setbox1\hbox{${\scriptscriptstyle #1}_{\scriptscriptstyle #2}$}
              \restrictionaux{#1}{#2}}}
\def\restrictionaux#1#2{{#1\,\smash{\vrule height .8\ht1 depth .85\dp1}}_{\,#2}}
%%%%%%%%%%%

%% SEZIONE GRAFICA
\use{tikz}
\usetikzlibrary{matrix, patterns, calc, decorations.pathreplacing, hobby, decorations.markings, decorations.pathmorphing, babel}
\use{tikz-3dplot}
\use{mathrsfs} %per geogebra
\use{tikz-cd}
\tikzset
{
  %surface/.style={fill=black!10, shading=ball,fill opacity=0.4},
  plane/.style={black,pattern=north east lines},
  curve/.style={black,line width=0.5mm},
  dritto/.style={decoration={markings,mark=at position 0.5 with {\arrow{Stealth}}}, postaction=decorate},
  rovescio/.style={decoration={markings,mark=at position 0.5 with {\arrow{Stealth[reversed]}}}, postaction=decorate}
}
\use{pgfplots} % stampare le funzioni
	\pgfplotsset{/pgf/number format/use comma,compat=1.15}
	%\pgfplotsset{compat=1.15} %per geogebra
	\usepgfplotslibrary{fillbetween, polar}
%%%%%%

%% CITAZIONI
\use{lineno}

% % Change standard block colors

\setbeamerfont{block title}{series=\bfseries} % Grassetto nei titoli

% \setbeamertemplate{blocks}[default]   %%  PRESENTAZIONE SQUADRATA
\setbeamertemplate{blocks}[rounded][shadow=true]

% % 1- Block title (background and text)

\setbeamercolor{block title}{use={palette primary,palette secondary}, bg=palette secondary.bg, fg=palette primary.fg}

% % 2- Block body (background)
\setbeamercolor{block body}{use={palette primary}, bg=palette primary.bg}

% % Change alert block colors

% % 1- Block title (background and text)
\setbeamertemplate{blocks}[rounded][shadow=true]
\setbeamercolor{block title alerted}{use={palette secondary}, bg=palette secondary.fg, fg=white}

% % 2- Block body (background)
\setbeamercolor{block body alerted}{use={palette primary}, bg=palette primary.bg}


% {\usebeamercolor[fg]{palette secondary} Prova} % colore giusto per lo sfondo
% {\usebeamercolor[bg]{palette secondary} Due} % colore giusto per i titoli

\newenvironment{definizioneb}{\begin{block}{Definizione}}{\end{block}}
\newenvironment{osservazioneb}{\begin{block}{Osservazione}}{\end{block}}
\newenvironment{teoremab}{\begin{alertblock}{Teorema}}{\end{alertblock}}
\newenvironment{corollariob}{\begin{alertblock}{Corollario}}{\end{alertblock}}

% \setbeamertemplate{itemize items}[square] %%  PRESENTAZIONE SQUADRATA
\setbeamertemplate{itemize items}[circle]
\setbeamercolor{itemize item}{use={palette primary}, fg=palette primary.fg}
\setbeamertemplate{enumerate items}[default]
\setbeamercolor{enumerate item}{use={palette primary}, fg=palette primary.fg}

\setbeamertemplate{navigation symbols}{}%remove navigation symbols

\newcommand{\citazione}[1]{%
  \begin{quotation}
  \begin{linenumbers}
  \modulolinenumbers[5]
  \begingroup
  \setlength{\parindent}{0cm}
  \noindent #1
  \endgroup
  \end{linenumbers}
  \end{quotation}\setcounter{linenumber}{1}
  }

\bibliography{bibliography.bib}

\hypersetup{%
	pdfauthor={Davide Peccioli},
	pdfsubject={},
	allcolors=black,
	citecolor=black,
	colorlinks=true,
	bookmarksopen=true}

\renewcommand{\href}[2]{#2}
\usetheme{CambridgeUS}
\usecolortheme{spruce}
\usefonttheme{professionalfonts}
\author{Davide Peccioli}
\date{17 agosto 2026}
\title{Teorema di Compattezza e Numeri Reali non-standard}
\begin{document}

\maketitle
\begin{frame}[label={sec:orgddb237d}]{Domanda}
Consideriamo i numeri reali \(\R\).

\begin{itemize}
\item Esiste un numero positivo \(\varepsilon\) che sia \alert{minore} di tutti gli altri numeri reali positivi?
\item La risposta è \alert{no}. Ma possiamo provare a costruire un insieme più grande di \(\R\), chiamato \(\null^{*}{\R}\), in cui questo è vero.
\end{itemize}

Questo è il nostro obiettivo, e per raggiungerlo avremo bisogno del \alert{Teorema di Compattezza}.
\end{frame}
\begin{frame}[label={sec:orgb21d814}]{Teorema di Compattezza}
\begin{block}{Teorema}
Sia \(T\) una teoria finitamente soddisfacibile. Allora \(T\) è soddisfacibile.
\end{block}
Per capire questo teorema, è necessario spiegare cosa siano:
\begin{itemize}
\item una teoria;
\item una teoria \emph{soddisfacibile};
\item una teoria \emph{finitamente} soddisfacibile.
\end{itemize}
\end{frame}
\begin{frame}[label={sec:org0085cfb}]{Partiamo dall'inizio}
\begin{block}{Esempio}
Consideriamo la frase
\begin{equation*}
P:\qquad \text{``Tra due numeri, ne esiste sempre uno in mezzo.''}
\end{equation*}
Questa frase \alert{è vera}?
\end{block}
La risposta è: \alert{dipende}.

Il compito della logica è proprio formalizzare la verità.
\end{frame}
\begin{frame}[label={sec:org85d2f9b}]{\null}
Ricordiamo che
\begin{align*}
\N &= \set{0,1,2,\dots,}\\
\Z &= \set{\dots, -3,-2,-1,0,1,2,\dots}\\
\Q &= \set{-\frac{1}{4}, \dots ,\frac{1}{2}, 0, 1, \frac{3}{2}, \dots}.
\end{align*}

Svisceriamo l'esempio di prima, e consideriamo \(P\) ``dentro'' \(\N, \Z\) e \(\Q\).
\begin{itemize}
\item In \(\N\) e \(\Z\), \(P\) è \alert{falsa}:
\begin{equation*}
  1\text{ e } 2,\qquad -1\text{ e } 0, \qquad \dots
\end{equation*}
\item In \(\Q\), \(P\) è \alert{vera}.
\end{itemize}
\end{frame}
\begin{frame}[label={sec:orga2d88ab}]{Modelli ed enunciati}
Questa cosa si riassume scrivendo che
\begin{equation*}
\N \nvDash P,\qquad \Z\nvDash P, \qquad \Q\vDash P
\end{equation*}
dove ``\(\vDash\)'' si legge ``modella''.

\begin{itemize}
\item \(\N, \Z\) e \(\Q\), che sono \alert{il luogo in cui ha senso parlare di verità}, si chiamano ``modelli'' o ``strutture''.
\item La \(P\) di prima, invece, si chiama \alert{enunciato}.

Un enuciato è una affermazione che può essere vera o falsa \alert{dentro ad un modello}.
\end{itemize}
\end{frame}
\begin{frame}[label={sec:org1c92fc1}]{Teoria}
\begin{block}{Definizione}
Una \alert{teoria} è un insieme di enunciati.
\end{block}
\begin{block}{Definizione}
Una \alert{teoria soddisfacibile} è un insieme di enunciati \alert{tutti veri dentro ad un modello \(\mathcal{M}\)}.
\end{block}
\begin{block}{Definizione}
Una teoria si dice \alert{finitamente soddisfacibile} se ogni suo sottoinsieme finito è soddisfacibile.
\end{block}
\end{frame}
\begin{frame}[label={sec:org6dd185a}]{Teorema di Compattezza}
A questo punto possiamo enunciare nuovamente il:
\begin{block}{Teorema}
Sia \(T\) una teoria finitamente soddisfacibile. Allora \(T\) è soddisfacibile.
\end{block}
Dimostrare questo teorema è un arduo compito, non adatto a questa serata. Però possiamo \alert{utilizzarlo} per costruire, ``a mano'', il nostro \(\null^{*} \R\).
\end{frame}
\begin{frame}[label={sec:orge137b2f}]{Teoria per i reali non-standard}
Consideriamo la teoria \(T\), composta da ``tutti'' gli \alert{enunciati veri in \(\R\)}.

Dentro \(T\) ci sono cose come:
\begin{itemize}
\item \(1+1=2\);
\item ``Ci sono infiniti numeri'';
\item ``Tra due numeri ne esiste sempre uno in mezzo'';
\item \ldots{}
\end{itemize}

Se consideriamo un simbolo in più, chiamato, ad esempio, \(\varepsilon\), aggiungiamo a \(T\) una quantità infinita di enunciati:
\begin{equation*}
0 < \varepsilon < 1, \quad 0 < \varepsilon < 0.0012324, \quad 0 < \varepsilon < \pi, \quad \dots
\end{equation*}
\end{frame}
\begin{frame}[label={sec:orgfb10bf5}]{La teoria è soddisfacibile?}
A questo punto, la teoria che abbiamo costruito, è finitamente soddisfacibile?

\alert{Sì}, e per ciascun sottoinsieme finito una struttura adatta è proprio \(\R\).
\end{frame}
\begin{frame}[label={sec:orgbaeda65}]{Concludendo\ldots{}}
È quindi possibile applicare il \alert{Teorema di Compattezza}: poiché la teoria è finitamente soddisfacibile, allora è anche \alert{soddisfacibile}.

Esiste quindi una struttura che rende vera \alert{tutti gli enunciati della teoria contemporaneamente}.

Questa struttura è proprio \(\null^{*} \R\). Quell'\(\varepsilon\) che abbiamo aggiunto sarà l'elemento positivo più piccolo di \alert{tutti gli altri elementi positivi}.
\end{frame}
\end{document}
